\section{De la graficación por computadora a la robótica: cómo el control en simulación llega al mundo real}

Esta sección comienza con la trayectoria de investigación de Tan Jie (谭杰). Su doctorado fue en graficación por computadora, en particular en physics-based character animation, es decir, lograr que figuras humanoides o animales caminen de forma natural en un entorno de simulación. Esta línea se parece mucho a la robótica: la graficación por computadora hace robótica en simulation, y la robótica hace graficación por computadora en el mundo real. La diferencia es que en la simulación se tiene acceso a toda la información de estado, mientras que el mundo real está lleno de ruido, fricción e incertidumbre.

\begin{figure}[H]
\centering
\includegraphics[width=0.96\textwidth]{graphics-to-robotics.png}
\caption{De la graficación por computadora a la robótica: hacer robótica en simulación y después transferirla al mundo real. Diagrama conceptual propio, elaborado a partir de la conversación de 00:02:00--00:13:06.}
\end{figure}

\begin{knowledgebox}{Lectura de la figura: la graficación por computadora aportó la ruta de «aprender primero en simulación»}
La figura parte de la graficación por computadora y la animación física, pasa por el aprendizaje de la marcha mediante aprendizaje por refuerzo y llega, mediante la transferencia Sim2Real, al robot real. Lo central del trabajo temprano de Tan Jie fue llevar la capacidad de control aprendida en simulación al movimiento real de robots cuadrúpedos.
\end{knowledgebox}

\subsection{Sim2Real y el aprendizaje por refuerzo: el primer cambio de paradigma}

Esta subsección aborda primero el problema de «cómo se mueve el cuerpo». Antes de la aparición de los modelos grandes, la robótica ya había vivido un giro importante: que el robot aprenda primero a moverse en simulación mediante aprendizaje por refuerzo y después transfiera lo aprendido al hardware real.

Sim2Real es la abreviatura de simulation-to-real: primero se entrena una política en un entorno de simulación y luego se transfiere al robot real. Tan Jie mencionó su trabajo temprano en Google, Sim2Real Learning Agile Locomotion for Quadruped Robots, que usó aprendizaje por refuerzo profundo para resolver la locomoción ágil de robots cuadrúpedos. El aprendizaje por refuerzo es un método de entrenamiento en el que un agente prueba y se equivoca dentro de un entorno y mejora su comportamiento a partir de la señal de reward; PPO es uno de los algoritmos más usados.

\begin{importantbox}{El primer cambio de paradigma}
Antes, el control de movimiento de los robots dependía en gran medida del control clásico y del Model Predictive Control; en los últimos cinco a diez años, el aprendizaje por refuerzo + Sim2Real transformó prácticamente el campo de la marcha y la locomotion, y volvió comunes con rapidez capacidades motrices como correr, saltar, dar volteretas o boxear.
\end{importantbox}

\subsection{El segundo cambio de paradigma, traído por los modelos grandes}

A continuación se aborda el problema de «si el robot entiende la tarea». El control de movimiento resuelve el cuerpo; los modelos grandes aportan el lenguaje, el sentido común y la planificación. Por eso Tan Jie compara los modelos grandes con el cerebro.

Antes de los modelos grandes, muchos robots carecían de common sense y tampoco entendían el lenguaje natural. Para que un robot preparara café, había que escribir un programa o descomponer la tarea en instrucciones de control muy específicas. El cambio que trajeron los modelos grandes es que el robot puede entender lenguaje natural, descomponer una tarea en pasos y saber, por sentido común, que «preparar café» requiere más o menos una taza, agua, café, calentar y servir. Tan Jie usa la analogía del cerebro y el cerebelo: el modelo grande se parece más al cerebro, encargado del pensamiento, el lenguaje y la planificación; el aprendizaje por refuerzo se parece más al cerebelo, encargado del equilibrio, el movimiento y la ejecución.

\begin{figure}[H]
\centering
\includegraphics[width=0.94\textwidth]{brain-cerebellum-robot.png}
\caption{Cerebro y cerebelo: el modelo multimodal se encarga de la cognición y el aprendizaje por refuerzo, de la ejecución. Diagrama conceptual propio, elaborado a partir de la conversación de 00:02:00--00:13:06.}
\end{figure}

\begin{knowledgebox}{Lectura de la figura: el robot necesita cerebro y también cerebelo}
El cerebro se encarga de la comprensión del lenguaje, el sentido común y la planificación de tareas; el cerebelo, de la marcha, el equilibrio, la prensión y el control de bajo nivel. Con solo cerebro, el robot «sabe qué hacer pero no lo hace bien»; con solo cerebelo, el robot «puede moverse pero no entiende el objetivo».
\end{knowledgebox}

\subsection{Resumen de la sección}

La trayectoria de investigación de Tan Jie explica los dos puntos de inflexión de la robótica en la última década: primero, el aprendizaje por refuerzo y Sim2Real transformaron el control de movimiento; segundo, los modelos grandes llevaron el lenguaje, el sentido común y la planificación a los robots. La investigación en robótica actual tiene que ocuparse a la vez del cerebro y del cerebelo.
